% ============================================================
\chapter*{Pr\'eface}
\addcontentsline{toc}{chapter}{Pr\'eface}
% ============================================================

\begin{center}
\textit{<<\,Tous les mod\`eles sont faux, mais certains sont utiles.\,>>}\\[4pt]
--- George E.\ P.\ Box
\end{center}

\bigskip

La mod\'elisation math\'ematique est l'art et la science de traduire des ph\'enom\`enes
du monde r\'eel en structures math\'ematiques rigoureuses, de les analyser pour en
d\'eduire des pr\'edictions, et de confronter ces pr\'edictions aux observations.
Ce cours s'adresse aux \'etudiants de deuxi\`eme et troisi\`eme ann\'ee de licence
(L2/L3) en math\'ematiques, physique, biologie quantitative ou sciences de
l'ing\'enieur, disposant d'une base solide en analyse, alg\`ebre lin\'eaire
et \'equations diff\'erentielles.

\section*{Objectifs p\'edagogiques}

\`A l'issue de ce cours, l'\'etudiant sera capable de :

\begin{enumerate}
  \item \textbf{Formuler} un probl\`eme concret sous forme d'un mod\`ele
        math\'ematique, en identifiant variables, param\`etres et hypoth\`eses.
  \item \textbf{Analyser} qualitativement et quantitativement les \'equations
        obtenues : points d'\'equilibre, stabilit\'e, bifurcations,
        comportement asymptotique.
  \item \textbf{Simuler} num\'eriquement les mod\`eles \`a l'aide de Python
        (\texttt{scipy}, \texttt{numpy}, \texttt{matplotlib}).
  \item \textbf{Valider} un mod\`ele en le confrontant aux donn\'ees
        exp\'erimentales et en discutant ses limites.
  \item \textbf{Communiquer} les r\'esultats de mani\`ere claire et
        rigoureuse, tant par \'ecrit que par des visualisations.
\end{enumerate}

\section*{Philosophie du cours}

Nous adoptons une d\'emarche \emph{du concret vers l'abstrait}. Chaque chapitre
part d'un ph\'enom\`ene observable --- croissance d'une population, propagation
d'une \'epid\'emie, diffusion de la chaleur, fluctuation d'un prix --- et
construit pas \`a pas le mod\`ele math\'ematique associ\'e.  Les outils th\'eoriques
(analyse dimensionnelle, th\'eorie de la stabilit\'e, processus stochastiques)
sont introduits \emph{au moment o\`u ils deviennent n\'ecessaires}, et non comme
des pr\'erequis abstraits.

\begin{attention}[title=\textbf{Mod\`ele $\neq$ R\'ealit\'e}]
Un mod\`ele est toujours une simplification.  La citation de Box en
exergue doit rester pr\'esente \`a l'esprit tout au long du cours.
Nous insisterons syst\'ematiquement sur les \textbf{hypoth\`eses} de chaque
mod\`ele et sur les cons\'equences de leur violation.
\end{attention}

\section*{Structure du cours}

Le cours est organis\'e en dix chapitres, suivant une progression logique :

\begin{description}
  \item[Chapitre 1 --- Le processus de mod\'elisation.]
    M\'ethodologie g\'en\'erale, \'etapes de la mod\'elisation, classification
    des mod\`eles.
  \item[Chapitre 2 --- Analyse dimensionnelle et mise \`a l'\'echelle.]
    Th\'eor\`eme de Buckingham~$\Pi$, adimensionnement, ordres de grandeur.
  \item[Chapitre 3 --- Dynamique des populations.]
    Mod\`eles de Malthus, Verhulst, Lotka--Volterra ; stabilit\'e et portraits
    de phase.
  \item[Chapitre 4 --- Mod\`eles \'epid\'emiologiques.]
    SIR, SEIR ; nombre de reproduction de base $\mathcal{R}_0$ ;
    strat\'egies de vaccination.
  \item[Chapitre 5 --- Syst\`emes dynamiques discrets et chaos.]
    Application logistique, exposants de Lyapunov, diagrammes de bifurcation.
  \item[Chapitre 6 --- Mod\`eles de diffusion et de transport.]
    \'Equation de la chaleur, advection-diffusion, loi de Fick.
  \item[Chapitre 7 --- Mod\`eles \'economiques et financiers.]
    Mod\`ele de Solow, Black--Scholes, \'equilibre offre-demande.
  \item[Chapitre 8 --- Optimisation et mod\`eles variationnels.]
    Programmation lin\'eaire, conditions KKT, calcul des variations.
  \item[Chapitre 9 --- Mod\`eles stochastiques.]
    Marches al\'eatoires, cha\^ines de Markov, mouvement brownien, EDS.
  \item[Chapitre 10 --- \'Etudes de cas et projets.]
    Applications int\'egratives couvrant plusieurs domaines.
\end{description}

\section*{Pr\'erequis}

\begin{itemize}
  \item \textbf{Analyse :} suites, s\'eries, calcul diff\'erentiel et int\'egral
        (une et plusieurs variables), \'equations diff\'erentielles ordinaires.
  \item \textbf{Alg\`ebre lin\'eaire :} espaces vectoriels, matrices,
        valeurs propres, diagonalisation.
  \item \textbf{Probabilit\'es :} espaces probabilis\'es, variables al\'eatoires,
        esp\'erance, variance, loi normale.
  \item \textbf{Informatique :} bases de Python ; les compl\'ements
        n\'ecessaires seront introduits au fil du cours.
\end{itemize}

\section*{Conventions et notations}

\begin{itemize}
  \item $\R$, $\N$, $\mathbb{Z}$, $\mathbb{Q}$, $\mathbb{C}$ : ensembles
        de nombres usuels.
  \item $\E[X]$, $\Var(X)$ : esp\'erance et variance d'une variable
        al\'eatoire $X$.
  \item $\norm{\cdot}$, $\abs{\cdot}$ : norme et valeur absolue.
  \item Vecteurs en gras : $\mathbf{x}$, $\mathbf{v}$.
  \item D\'eriv\'ees temporelles : $\dot{x} = dx/dt$.
  \item Grandeurs adimensionn\'ees : $\tilde{x} = x/x_0$.
\end{itemize}

\section*{Organisation pratique}

Chaque chapitre contient :
\begin{itemize}
  \item Des \textbf{d\'efinitions}, \textbf{th\'eor\`emes} et
        \textbf{propositions} num\'erot\'es.
  \item Des \textbf{exemples d\'etaill\'es} illustrant les concepts.
  \item Des encadr\'es \textbf{Formules cl\'es}, \textbf{Attention},
        \textbf{Bonne pratique} et \textbf{Intuition}.
  \item Des \textbf{simulations Python} compl\`etes et comment\'ees.
  \item Des \textbf{exercices} de difficult\'e croissante.
\end{itemize}

\section*{Logiciels}

\begin{codeblock}[title=\textbf{Biblioth\`eques Python utilis\'ees}]
\begin{minted}{python}
import numpy as np
import matplotlib.pyplot as plt
from scipy.integrate import odeint, solve_ivp
from scipy.optimize import minimize, linprog
from scipy.linalg import eig
\end{minted}
\end{codeblock}

\section*{Remerciements}

Ce cours s'inspire des ouvrages de C.~Lin et L.~Segel
(\textit{Mathematics Applied to Deterministic Problems in the Natural Sciences}),
J.\,D.~Murray (\textit{Mathematical Biology}),
S.\,H.~Strogatz (\textit{Nonlinear Dynamics and Chaos}) et
E.\,A.~Bender (\textit{An Introduction to Mathematical Modeling}).

\bigskip
\begin{flushright}
\textit{Bonne lecture et bonne mod\'elisation\,!}
\end{flushright}
